% TREŚĆ ROZDZIAŁU: zastąp wskazówki opisem wykonania rozwiązania.
% Nie wklejaj długiego kodu bezpośrednio do pliku .tex. Umieść go w katalogu
% code/ rozdziału i wczytaj przez \lstinputlisting — przykład jest na końcu.
\chapter{Realizacja projektu}
\label{chap:realizacja}

\section{Środowisko i narzędzia}
Wymień tylko te wersje narzędzi, bibliotek, urządzeń i konfiguracji, które są istotne dla odtworzenia projektu.

\section{Implementacja lub wykonanie}
Funkcja \texttt{max\_element} zwraca największy element niepustej tablicy liczb całkowitych. Jej kod przedstawia \cref{lst:max-element}.

\lstinputlisting[
  style=project-c,
  caption={Wyszukiwanie największego elementu tablicy},
  label={lst:max-element}
]{chapters/04-realizacja/code/max_element.c}

Składnię języka C i działanie kompilatora opisuje dokumentacja GCC \cite{gcc2026}.
Listing wczytano mechanizmem pakietu \texttt{listings}, opisanym w dokumentacji LaTeX-a \cite{latexdocs2026}.
\subsection{Implementacja algorytmu}

Zamianę sąsiednich elementów zapisano w \cref{lst:porownanie}.

\begin{lstlisting}[
  style=project-python,
  caption={Porównanie dwóch sąsiednich elementów},
  label={lst:porownanie}
]
if result[j] > result[j + 1]:
    result[j], result[j + 1] = result[j + 1], result[j]
\end{lstlisting}

Kompletny program pochodzi z osobnego pliku i jest pokazany
w \cref{lst:bubble-sort}.

\lstinputlisting[
  style=project-python,
  caption={Sortowanie bąbelkowe},
  label={lst:bubble-sort}
]{chapters/04-realizacja/code/bubble_sort.py}

\section{Instrukcja uruchomienia}
Podaj minimalny, jednoznaczny zestaw kroków potrzebnych do uruchomienia, zbudowania albo użycia rozwiązania.

% Kod przechowuj w katalogu code/ odpowiedniego rozdziału i wczytuj np.:
% \lstinputlisting[style=project-python,caption={Opis kodu},label={lst:przyklad}]{chapters/04-realizacja/code/program.py}
